% ============================================================
% portada.tex — Título, candidato, director(es), resumen y abstract
% ============================================================

\begin{titlepage}
  \centering
  \vspace*{0.4cm}

  \includegraphics[width=3.4cm]{figuras/logo-utm.png}\par
  \vspace{0.5cm}

  {\scshape\Large Universidad Tecnológica de la Mixteca}\par
  \vspace{2pt}
  {\large Maestría en Ingeniería de Software}\par

  \vspace{2cm}

  {\rule{0.55\textwidth}{0.6pt}}\par
  \vspace{0.4cm}
  {\Large \scshape Protocolo de Tesis de Maestría}\par
  \vspace{0.4cm}
  {\rule{0.55\textwidth}{0.6pt}}\par

  \vspace{1.8cm}

  {\large \textit{Título:}}\par
  \vspace{0.4cm}
  {\Large \bfseries Evaluación de un modelo de orquestación
  multiagente con roles especializados y heterogeneidad de modelos
  frente a un agente único en tareas de desarrollo de software \par}

  \vspace{2cm}

  \begin{tabular}{@{}r@{\hspace{0.6em}}l@{}}
    \textbf{Candidato:} & Luis David Pérez Cruz \\[0.5cm]
    \textbf{Director(a) de Tesis:} & Carlos Alberto Fernández y Fernández \\
    %\textbf{Codirector(a) de Tesis:} & Nombre del codirector (si aplica) \\
  \end{tabular}

  \vfill
  {\rule{0.35\textwidth}{0.4pt}}\par
  \vspace{0.3cm}
  {\large \the\year \par}
\end{titlepage}

\section*{Resumen}
El desarrollo de software asistido por IA avanza hacia agentes de codificación que ejecutan tareas de forma autónoma, lo que ha generado dudas sobre la corrección y la estructura del código que producen. El Role Specialization Model (RSM) propone coordinar varias herramientas basadas en modelos de lenguaje asignándoles roles especializados, pero carece de evaluación empírica y cuantitativa. Este trabajo evaluará si un modelo de orquestación multiagente con roles especializados y heterogeneidad completa de modelos (cada rol ocupado por un modelo de una familia o proveedor distinto), basado en el RSM y operado por un orquestador humano, produce código con mayor corrección funcional y mejor estructura arquitectónica que un agente único de propósito general. Para ello, se diseñará e implementará el artefacto (el modelo de orquestación) y se evaluará mediante un experimento controlado de diseño pareado sobre un banco de tareas de desarrollo de backend en TypeScript. La corrección funcional se medirá por la compilabilidad y el \textit{pass rate}, y la estructura por la modularidad y la complejidad cognitiva, conforme a la norma ISO/IEC 25010. La comparación entre ambos enfoques se realizará con pruebas estadísticas pareadas. Se espera determinar si la combinación de roles especializados y heterogeneidad de modelos mejora la corrección funcional y la estructura del código frente a un agente único.

\section*{Palabras clave}
orquestación multiagente, roles especializados, heterogeneidad de modelos, desarrollo de software asistido por IA, corrección funcional del código

\newpage
